\chapter{Benchmark Environment}
\label{ch:benchmark}

\section{Why Competitive Pok\'emon}

Competitive Pok\'emon Singles (CPS) is a turn-based strategy game with
simultaneous move selection each turn. Players choose teams of six Pok\'emon
with moves, items, abilities, and in Generation 9 a terastallization type. At
battle start much of the opponent roster is hidden. Species, sets, and intent
must be inferred from damage ranges, switch patterns, item triggers, and move
reveals.

CPS therefore combines partial observability, simultaneous commitment, long
horizons (often 20--100+ turns), and combinatorial team structure. Human replay
data spans more than a decade and grows daily on public ladders.

\section{Infrastructure}

We use three open components.
\begin{itemize}
  \item \textbf{poke-env} bridges Python agents to a local Pokemon Showdown
    server for live evaluation.
  \item \textbf{Metamon} reconstructs first-person trajectories from spectator
    replays with full-state annotations for belief labels \citep{grigsby2023metamon}.
  \item \textbf{Gymnasium} provides the environment interface standard for RL
    tooling in this project.
\end{itemize}

\section{Scope of This Chapter}

This thesis is not about Pok\'emon mechanics, tier lists, or heuristic engines.
We use Gen 9 OU human battles as the primary corpus. Random-bot logs from
poke-env support pipeline debugging only. All scientific claims concern belief
inference and generalization, with CPS as an empirical testbed.

\section{Action Space}

At each decision point the agent may use a move targeting an opponent slot or
switch to a bench Pok\'emon. Actions are encoded as discrete classes in an
action vocabulary built from training data. The exact action cardinality depends
on format rules and logged move names.

\section{Data Pipeline Status}

The implementation roadmap proceeds in phases.
\begin{enumerate}
  \item Local Showdown server and baseline logging \textbf{(complete)}.
  \item 13-token turn encoder and vocabulary \textbf{(complete)}.
  \item TurnTransformer behavioral cloning \textbf{(complete)}.
  \item Metamon download and bridge to token format \textbf{(in progress)}.
  \item Belief label extraction and multi-task training \textbf{(in progress)}.
  \item Combinatorial split tooling and evaluation harness \textbf{(planned)}.
\end{enumerate}

\Cref{app:implementation} lists repository paths for each component.
